\section{Problem 4: Program Verification}
A C program uses nondeterministic choices to update two variables before checking an assertion.
Determine which values of $k$ can cause the program to crash and provide an execution for each possible crash.

\subsection{Solution}
We identify multiple parameters for the problem
\begin{itemize}
\item $q$ the iteration bound. We set this to the maximum amount the loop would be able to iterate which is 40.
\item $a_i \in \mathbb{N}$ the value of the $a$ variable of the $i$th iteration
\item $b_i \in \mathbb{N}$ the value of the $b$ variable of the $i$th iteration
\item $k \in \{0..10\}$ the constant added to the last \verb|if| statement.
\end{itemize}

The begin of the snippet defines $a$ and $b$ to 1. Which we encode as 
\[
a_0 = 1 \wedge b_0 = 1
\]

Encoding each iterations we have bound it to $q$. Leading to the formula.
\begin{equation} \label{eq:e4-2}
\bigwedge_{1 \le i \le q} F
\end{equation}
where $F$ is the formula of the iteration body.

Since $q$ is the maximum number of iterations not all runs need 40 iterations. Therefore we guard the formula with the condition of the while loop:
\begin{equation}
\big( \  a_{i-1} < 180 \Rightarrow F) \big) \wedge 
\big( \ a_{i-1} \geq 180 \Rightarrow ( b_i = b_{i-1} \wedge a_i = a_{i-1} )\big)
\end{equation}
where $F$ is the formula of the iteration body.

Encoding the case \verb|__nondet__()| returns true we get
\[
(b_i = b_{i-1} + 3 \wedge a_i = a_{i-1} + 2 * b_i )
\]
Encoding the case where \verb|__nondet__()| returns false we get
\[
(b_i = b_{i-1} + a_{i-1} \wedge a_i = a_{i-1} + 5 )
\]

Encoding a \verb|if-then-else| statement normally would result in the form $(c \Rightarrow F_1) \wedge (\neg c \Rightarrow F_2)$.
Since the condition is non deterministic we cannot encode the condition that way, we treat it as either case is correct, therefore we encode it as $F_1 \vee F_2$.
\begin{equation} \label{eq:e4-1}
(b_i = b_{i-1} + 3 \wedge a_i = a_{i-1} + 2 * b_i ) \vee
(b_i = b_{i-1} + a_{i-1} \wedge a_i = a_{i-1} + 5 )
\end{equation}


Normally for program correctness you would prove \texttt{UNSAT} for $\neg Q$.
But in this case we are trying to verify than no trace could lead to the condition resulting in true.
So we would prove \texttt{UNSAT} for $Q$. This results in.
\[
b_q = 190 + k
\]

Combining all the components together we get the formula:
\begin{equation}
\begin{split}
	& a_0 = 1 \wedge b_0 = 1 \\
	& \bigg(\bigwedge_{1 \le i \le q}  \\
		& \quad \big( \  a_{i-1} < 180 \Rightarrow 
			( b_i = b_{i-1} + 3 \wedge a_i = a_{i-1} + 2 * b_i ) \vee
			( b_i = b_{i-1} + a_{i-1} \wedge a_i = a_{i-1} + 5 )
		\big) \wedge \\ 
		& \quad \big( \ a_{i-1} \geq 180 \Rightarrow 
			( b_i = b_{i-1} \wedge a_i = a_{i-1} )
		\big) \\
	& \bigg) \wedge b_q = 190 + k
\end{split}
\end{equation}

The program builds the formula and for each value for $k$ it checks the formula.

\subsection{Results}
For each K the program returns \texttt{SAT}, meaning the program can crash, or \texttt{UNSAT}, meaning the program cannot crash. It takes about 12 seconds to compute all the results.
\begin{table}[H]
\centering
\caption{Result of Satisfiability}
\label{tab:exercise4}
\begin{tabular}{c|l}
	\textbf{K} & \textbf{Result} \\
    \hline
    %%% T | N | P | S | C | D | W
	0  & Program can not crash.  \\
	1  & Program can crash.  \\
	2  & Program can not crash.  \\
	3  & Program can crash.  \\
	4  & Program can not crash.  \\
	5  & Program can crash.  \\
	6  & Program can crash.  \\
	7  & Program can crash.  \\
	8  & Program can not crash.  \\
	9  & Program can not crash.  \\
	10 & Program can not crash.  \\
\end{tabular}
\end{table}
Per \texttt{k} that can crash we show the values of \texttt{a}, \texttt{b}  and what \verb|__nondet__()| returned.
The trace is given until the loop would exit.
The formula would carry the value over across the rest of the iterations. These traces can be seen in the Table~\ref{tab:exercise4-traces}


\begin{table}[htbp]

\begingroup
	\newcommand\header[1]{
	\hline
		\tabletitle{Trace for k = #1} \\
	\hline
		Iter. & $a$  & $b$ & Condition \\
	\hline}
\caption{Traces for exercise 4}
\label{tab:exercise4-traces}
\def\false{ \texttt{False} }
\def\true{\text{True}}
	\def\tabletitle #1{\multicolumn{4}{c}{\cellcolor{black}\color{white}#1}}
\centering
\begin{tabular}{ccc}

	\begin{tabular}[t]{|rrrr|}
	\header{1}
		0 	  & 1         & 1         & - \\
		1 	  & 6         & 2         & \false \\
		2 	  & 16        & 5         & \true \\
		3 	  & 32        & 8         & \true \\
		4 	  & 54        & 11        & \true \\
		5 	  & 59        & 65        & \false\\
		6 	  & 64        & 124       & \false \\
		7 	  & 69        & 188       & \false \\
		8 	  & 451       & 191       & \true \\
		9 	  & 451       & 191       & - \\
		\multicolumn{4}{|c|}{Left out for brevity} \\
	\hline
	\end{tabular}
	&
	\begin{tabular}[t]{|rrrr|}
	\header{3}
		0  	  & 1    & 1   & -\\
		1  	  & 6    & 2   & \false \\
		2  	  & 11   & 8   & \false \\
		3  	  & 16   & 19  & \false \\
		4  	  & 21   & 35  & \false \\
		5  	  & 26   & 56  & \false \\
		6  	  & 31   & 82  & \false \\
		7  	  & 36   & 113 & \false \\
		8  	  & 41   & 149 & \false \\
		9  	  & 46   & 190 & \false \\
		10 	  & 432  & 193 & \true \\
		11 	  & 432  & 193 & - \\
		\multicolumn{4}{|c|}{Left out for brevity} \\
	\hline
	\end{tabular}
	&
	\begin{tabular}[t]{|rrrr|}
	\header{5}
		0  	  & 1   & 1    & -\\
		1  	  & 9   & 4    & \true \\
		2  	  & 14  & 13   & \false \\
		3  	  & 46  & 16   & \true \\
		4  	  & 84  & 19   & \true \\
		5  	  & 89  & 103  & \false \\
		6  	  & 94  & 192  & \false \\
		7  	  & 484 & 195  & \true \\
		8  	  & 484 & 195  & - \\
		\multicolumn{4}{|c|}{Left out for brevity} \\
	\hline
	\end{tabular}
	\\ 
	\begin{tabular}[t]{|rrrr|}
	\header{6}
		0  	  & 1   & 1    & -\\
		1  	  & 9   & 4    & \true \\
		2  	  & 14  & 13   & \false \\
		3  	  & 19  & 27   & \false \\
		4  	  & 79  & 30   & \true \\
		5  	  & 84  & 109  & \false \\
		6  	  & 89  & 193  & \false \\
		7  	  & 481 & 196  & \true \\
		8  	  & 481 & 196  & - \\
		\multicolumn{4}{|c|}{Left out for brevity} \\
	\hline
	\end{tabular}

	&

	\begin{tabular}[t]{|rrrr|}
	\header{7}
		0  	  & 1 & 1      & -\\
		1  	  & 6 & 2      & \false \\
		2  	  & 11 & 8     & \false \\
		3  	  & 16 & 19    & \false \\
		4  	  & 60 & 22    & \true \\
		5  	  & 110 & 25   & \true \\
		6  	  & 166 & 28   & \true \\
		7  	  & 171 & 194  & \false \\
		8  	  & 565 & 197  & \true \\
		9  	  & 565 & 197  & - \\
		\multicolumn{4}{|c|}{Left out for brevity} \\
	\hline
	\end{tabular}
	&\\
\end{tabular}
\endgroup
\end{table}

\subsection{Remarks}
For the formula we have bound $q$ to be 40, which is the maximum amount of iterations the loop would be able to make. This does lead to the formula to maybe bigger than it needs.
